\section{OctoML：从开源到创业}

\subsection{创业动机}

OctoML 的创立非常自然：TVM 项目有几个同学和老师，觉得有商业化机会，就决定做 startup。陈天奇的核心动机是\textbf{``希望开源社区能够成功''}------为了开源去做 startup，而不是为了 startup 去做开源。

他第一年全职参与，之后转为 part-time，主要负责技术方向和下一代技术研究。他``很明确不要当 CEO''，因为``你要放弃技术很多''，而他的 passion 在于做好的技术。

\subsection{商业模式的转变}

OctoML 的商业模式经历了显著转变：

\begin{enumerate}
    \item \textbf{早期（2019--2022）}：通过机器学习编译帮助硬件厂商和客户更快部署模型到各种硬件。潜在客户包括芯片厂商和需要部署 ML 的企业。
    \item \textbf{后期（2023--2024）}：转型为 LLM API 服务商（类似 Fireworks、Together），直接提供 Llama 等模型的 endpoint，靠 API 调用收费。机器学习编译技术退居幕后成为支撑技术，帮助降低 inference 成本。
    \item \textbf{2024 年底}：被 NVIDIA 收购，团队负责 NVIDIA 的 AI compiler 方向。
\end{enumerate}

主持人追问 API 市场竞争激烈，大家都在比谁的 Llama inference cost 更低，OctoML 如何应对？陈天奇承认竞争激烈，但指出产品形态上仍有差异化空间，比如支持 multi-LoRA inference 等。他也坦言，NVIDIA 收购并非预设的退出路径------``我完全没有想过这个问题，当时就觉得能够把一件事情做好看看就好。''

\begin{knowledgebox}{导师创业经历的预示}
在加入 OctoML 之前，陈天奇已经目睹了导师 Carlos Guestrin 的创业公司 Turi（前身 GraphLab/Dato）的演变：从 CMU 的大规模图处理软件，到被收购时已完全转型为 Data Science 平台。``所以在开始做这件事情的时候，就会有心理准备------你去开始做的事情和 what you end up doing 一定会是不一样的。'' Databricks 从 Spark 出发，现在 Spark 也不再是其最主要的部分。这是开源 startup 的普遍规律。
\end{knowledgebox}

\begin{warningbox}{开源与商业的微妙关系}
陈天奇坦承开源项目做 startup 会面临``开源太好用导致没人付费''的问题。但他的观点是：（1）开源是你的起点，但几乎一定不是终点------Databricks 现在的核心业务也不再是 Spark；（2）所有 startup 开始想做的东西和最后做的东西``一定是不一样的''；（3）既然选择了开源，就要接受它的好处（社区反馈、影响力）和挑战（商业模式设计）。
\end{warningbox}

\subsection{创业最大的 Lesson}

创业五年给陈天奇带来了技术之外的重要认知：

\begin{itemize}
    \item \textbf{Human factor is very important}：``做 startup 最大的 lesson 是------很多时候它不光有技术的因素。怎么样跟人交流，怎么样在大的团队中扮演一个 leadership 的 role，我觉得这一点非常重要。''
    \item \textbf{必须 reinvent yourself}：不能抱着原来的技术路线不变。``如果我们两年前还在继续优化 ResNet，这个我觉得最后我们会输得很惨。''在快速变化的 AI 领域，技术转型的速度和勇气至关重要。
    \item \textbf{技术 excellence 始终没有错}：``这个观点不会被所有人认同。但在我看来，很多时候做优秀的技术始终是没有错的。因为 in the end, you have to have something to differentiate yourself。''
    \item \textbf{理想主义 + 现实主义}：``我有一点点理想主义。但这段经历会让我觉得我会更现实------我会知道说怎么样能够通过现实的方式去更好地实现理想。''
    \item \textbf{Product Market Fit}：需要思考你的产品到底解决了谁的什么问题。但同时不能只看短期 PMF 而忽视技术的长期价值。
    \item \textbf{明确自己的角色}：``我很明确我不要当 CEO。因为当 CEO 很难很难，你要放弃技术很多。我的 passion 在于做好的技术。''了解自己擅长什么、热爱什么，然后找到合适的 co-founder 互补。
\end{itemize}

\begin{warningbox}{开源与商业的微妙关系}
陈天奇坦承开源项目做 startup 会面临``开源太好用导致没人付费''的问题（主持人举了 Anyscale 和 Ray 的例子）。但他的观点是：（1）开源是你的起点，但几乎一定不是终点------Databricks 现在的核心业务也不再是 Spark；（2）所有 startup 开始想做的东西和最后做的东西``一定是不一样的''；（3）既然选择了开源，就要接受它的好处（社区反馈、影响力、用户信任）和挑战（商业模式需要差异化）。上一期嘉宾胡渊鸣的太极（Taichi）也遇到了类似问题------``很多人都说太极很好，但不愿意花钱。''这是基于开源创业的普遍挑战。
\end{warningbox}

\subsection{本章小结}

OctoML 的五年历程展示了学术开源项目商业化的典型路径：从技术驱动到产品驱动，从编译工具到 API 服务，最终被战略收购。陈天奇在这个过程中学到了技术之外的重要维度：人的因素、产品市场匹配、以及持续自我革新的必要性。

